\chapter{Rezultati}\label{cha:rezultati}
V tem poglavju študent predstavi rezultate meritev, analiz, preračunov in drugih obdelav podatkov, pridobljenih v okviru raziskave. Rezultati naj bodo predstavljeni objektivno.

Običajno je delo obsežnejše in sestavljeno iz več ločenih sklopov, zato najpogosteje rezultate iz posameznega sklopa predstavljate v ločenih (pod)poglavjih. Končna oblika mora biti takšna, da je karseda pregledna, jasna in razumljiva.

Poglavji \ref{cha:rezultati} in \ref{cha:diskusija} lahko združite v enotno poglavje \ref{cha:rezultati} Rezultati in diskusija, v kolikor je za izboljšanje jasnosti in utemeljevanja raziskovalnih hipotez/ciljev to bolj primerno.
% Po potrebi naslov spremenite v »Rezultati in diskusija« in v glavni
% datoteki odstranite \chapter{Diskusija}\label{cha:diskusija}
V tem delu predstavite vaše razumevanje rezultatov, njihovo interpretacijo in morebitne komentarje. Diskusija naj se vseskozi karseda neposredno navezuje na poglavja, ki predstavljajo odprta vprašanja in obravnavano problematiko identificirana v poglavju \ref{cha:pregled} in predstavljene hipoteze/cilje v poglavju \ref{cha:hipoteze}.

Poglavji \ref{cha:rezultati} in \ref{cha:diskusija} lahko združite v enotno poglavje \ref{cha:rezultati} Rezultati in diskusija, v kolikor je za izboljšanje jasnosti in utemeljevanja raziskovalnih hipotez/ciljev to bolj primerno.
.
